\section{The outer tail and the half-electron radius}\label{a:sec:tail}

We now combine the positive ionization gap with a weighted shell
estimate.  The conditional field is estimated after changing the law
by the shell mass itself; the measurability of this mass in the cut
data is what avoids an additional particle-count factor.

\begin{proposition}[Uniform exterior mass]\label{a:prop:tail-mass}
There exist universal $C,R_0,Z_3<\infty$ such that every normalized
neutral ground state with $Z\ge Z_3$ satisfies
\begin{equation}\label{a:eq:tail-mass}
 \int_{|x|>R}\rho_{\Psi_Z}(x)\,\dd x\le\frac C R
 \qquad(R\ge R_0).
\end{equation}
\end{proposition}

\begin{proof}
Fix any such ground state $\Psi=\Psi_Z$ and write $\E$ for its raw
configuration expectation.  By Proposition~\ref{a:prop:tail-offset} its
offset is bounded by a universal constant, and by
Proposition~\ref{a:prop:tail-gap} its ionization gap is at least $c>0$.
For $u\ge1$, choose a smooth radial cutoff $\chi_u(r)=\chi(r/u)$,
where $0\le\chi\le1$, $\chi=1$ on $[1,2]$, and $\chi$ is supported
in $[1/2,4]$.  Put
\[
 \begin{gathered}
 w_i=\chi_u(|x_i|)^2,\qquad
 S_u=\sum_{i=1}^Z w_i,\qquad m_u=\E S_u,\\
 n_{\rm adj}=\#\{i:u/2\le|x_i|\le4u\},\qquad
 N_u=\E n_{\rm adj}.
 \end{gathered}
\]
In particular $m_u\le N_u$.  Define
\[
 \mathcal B_u(x)=V(x)-\sum_{|x_j|<u/8}K(x-x_j)
        =B_{u/8}(x).
\]
Only its values on the support of $\chi_u$ will be used, so the
selected electron never occurs among these inner sources.

Testing the weak eigen-equation by $S_u\Psi$, and decomposing the
result into its $i$-th weighted terms, gives a useful lower bound.
For fixed $x_i$ and spin, the other $Z-1$ variables are antisymmetric;
their energy is at least $E_{Z-1}$ times their squared norm.  The
selected electron's kinetic pairing is
\[
 \operatorname{Re}\frac12\int
 \nabla_i\overline\Psi\cdot\nabla_i(w_i\Psi)
 =\frac12\int w_i|\nabla_i\Psi|^2
       -\frac14\int(\Delta_iw_i)|\Psi|^2
 \ge-Cu^{-2}\E\1_{\{u/2\le|x_i|\le4u\}}.
\]
The integration by parts follows first for smooth form approximants;
$w_i$ and its first two derivatives are bounded, so it passes to the
form domain.  Keeping only the selected electron's repulsion from
inner sources, and discarding its other nonnegative pair terms,
therefore gives
\begin{equation}\label{a:eq:tail-weighted}
 c m_u\le C u^{-2}N_u
       +\E\sum_i w_i \mathcal B_u(x_i).
\end{equation}

Suppose first that $m_u>0$.  The normalized antisymmetric state
\[
 \Psi_u^*=\sqrt{S_u/m_u}\,\Psi
\]
is form admissible.  Indeed $\sqrt{S_u}$ is the Euclidean norm of the
vector $(\chi_u(|x_i|))_i$, hence is Lipschitz.  Where $S_u>0$,
\[
 |\nabla\sqrt{S_u}|^2
 =\frac{\sum_i\chi_u(|x_i|)^2|\nabla_i\chi_u(|x_i|)|^2}{S_u}
 \le C u^{-2};
\]
its gradient vanishes almost everywhere on its zero set.  In
particular this is bounded by $Cu^{-2}n_{\rm adj}$.  Identity
\eqref{eq:sector-multiplier} shows that its offset satisfies
\begin{equation}\label{a:eq:tail-biased-offset}
 \delta_u\le C+Cu^{-2}N_u/m_u<\infty.
\end{equation}
There is no assumed positive lower bound on $m_u$.

We reveal every position carrying a positive shell weight so that
$S_u$ is determined by the data. Use an auxiliary cut which is full
out on $u/2\le|x|\le4u$, supported
in $u/4<|x|<8u$, and has squared gradient cost at most $Cu^{-2}$.
Let $\omega$ denote the raw configuration, let $\ell$ denote these
cut labels, and let $Y=Y(\omega,\ell)$ be all recorded out data.  If
$Q(\dd\ell\mid\omega)$ is the conditional label kernel, use exactly
this kernel for both states.  Their joint laws are then
\begin{equation}\label{a:eq:tail-joint-bias}
 \begin{split}
 P(\dd\omega,\dd\ell)
  &=|\Psi(\omega)|^2\,\dd\omega\,Q(\dd\ell\mid\omega),\\
 P^*(\dd\omega,\dd\ell)
  &=\frac{S_u(\omega)}{m_u}\,P(\dd\omega,\dd\ell).
 \end{split}
\end{equation}
Integration in $\omega$ includes the spin sum.  Every particle with
positive $w_i$ is recorded by this cut, so
\[
 S_u(\omega)=s(Y):=
 \sum_{i\ {\rm recorded}}\chi_u(|x_i|)^2.
\]
This equality includes the random labels in the transition regions:
particles there have weight zero.  It follows from
\eqref{a:eq:tail-joint-bias} that
\begin{equation}\label{a:eq:tail-conditional-bias}
 \frac{\dd P_Y^*}{\dd P_Y}=\frac{s(Y)}{m_u},\qquad
 \E_{P^*}[F\mid Y]=\E_P[F\mid Y]\quad\text{on }\{s>0\}
\end{equation}
for every bounded measurable $F$, and then for integrable $F$ by
truncation wherever the conditional expectations are defined.  For
example, multiply the asserted equality by $s(Y)$ and integrate over
an arbitrary data event to check it directly.  At the slice level the
same fact says that $\sqrt{S_u/m_u}$ is constant in the surviving core
variables and cancels when their wavefunction is normalized.  The set
$s=0$ has zero $P^*$-mass and contributes nothing to the original
weighted sum.

Let
\[
 h_Y(x)=\E_P[\mathcal B_u(x)\mid Y],\qquad
 M(Y)=\sup_{u/2\le|x|\le4u}(h_Y(x))_+.
\]
On the broad annulus $u/4<|x|<6u$ all the inner sources are separated
from the evaluation point.  Lemma~\ref{a:lem:tail-harmonic} supplies
continuous harmonic versions of these fields.  Applying
\eqref{a:eq:tail-conditional-bias} first at a countable dense set of
points and then using continuity makes the conditional fields for
$P$ and $P^*$ equal on the entire broad annulus, outside one data-null
set on $\{s>0\}$.  The two data laws are equivalent on that set.
Thus the equality applies at recorded random positions as well as to
the positive supremum.

Use \eqref{a:eq:tail-distant-field} in the state $\Psi_u^*$.  It gives
\begin{equation}\label{a:eq:tail-biased-field}
 \|M\|_{L^2(P^*)}
 \le C\big(u^{-1}+\sqrt{\delta_u}\,u^{-1/2}\big).
\end{equation}
This is an application of the screening estimate for arbitrary finite
offsets, not of its small-offset refinement.  In particular the cut
condition is uniform even if $m_u$ is very small: as checked in
Lemma~\ref{a:lem:tail-harmonic}, its cell coefficients obey
$D_z^2/(a_zm_z)\le C$ for every $\delta_u$, because $a_zm_z\ge1$.

To see explicitly how the weighted count is used, introduce the
data-measurable measure
\[
 \eta_Y=\sum_{i\ {\rm recorded}}w_i\delta_{x_i},
 \qquad \eta_Y(\R^3)=s(Y).
\]
Conditional integration and \eqref{a:eq:tail-joint-bias} now yield
\begin{align}
 \E\sum_iw_i\mathcal B_u(x_i)
 &=\E_P\int h_Y\,\dd\eta_Y\notag\\
 &=m_u\E_{P^*}
      \left[\frac1{s(Y)}\int h_Y\,\dd\eta_Y\right]\notag\\
 &\le m_u\E_{P^*}M
 \le C m_u\big(u^{-1}+\sqrt{\delta_u}\,u^{-1/2}\big).
 \label{a:eq:tail-weighted-field}
\end{align}
The quotient may be defined as zero on $s=0$.  This exact change of
measure accounts for all shell weights, so there is no further
particle-count factor in the last bound.

Substituting \eqref{a:eq:tail-biased-offset} into
\eqref{a:eq:tail-weighted-field}, and then into
\eqref{a:eq:tail-weighted}, gives
\begin{align*}
 c m_u
 &\le Cu^{-2}N_u
    +Cm_u(u^{-1}+u^{-1/2})
    +Cu^{-3/2}\sqrt{m_uN_u}\\
 &\le C(u^{-2}+u^{-3/2})N_u
    +Cm_u(u^{-1}+u^{-1/2}).
\end{align*}
The second line uses $m_u\le N_u$.  When $m_u=0$ it holds trivially
without introducing a biased law.  For all sufficiently large $u$,
absorb the last term into the left side.  Applying
\eqref{a:eq:screen-annular} in the original state, whose offset is
uniformly bounded, gives $N_u\le C(1+\sqrt u)$ and hence
\begin{equation}\label{a:eq:tail-shell-mass}
 \E\#\{i:u\le|x_i|\le2u\}\le m_u\le C/u.
\end{equation}
Finally, sum over $u=2^jR$, $j\ge0$.  The density gives zero mass to
the intervening spheres, and $\sum_{j\ge0}(2^jR)^{-1}=2/R$.
This proves \eqref{a:eq:tail-mass}.
\end{proof}

\begin{proof}[Completion of the proof of Theorem~\ref{thm:radius}]
For all sufficiently large $Z$, \eqref{a:eq:tail-lower} supplies a fixed
$s_*>0$ for which the exterior mass exceeds $1/2$.  Monotonicity of
that mass as a function of the radius implies
$R(\Psi_Z)\ge s_*$.  Proposition~\ref{a:prop:tail-mass} supplies a fixed
large $R_*$ for which the exterior mass is at most $1/2$, and hence
$R(\Psi_Z)\le R_*$.

For clarity we give uniform arguments over the entire ground-state
space for each of the finitely many remaining charges.  Fix such a
$Z$.  The strict gap $E_Z<E_{Z-1}$ and the compactness conclusion in
Lemma~\ref{lem:atomic-binding} show that every sequence of normalized
neutral ground states has a strongly $L^2$ convergent subsequence.
In particular their configuration laws are uniformly tight.  To see
directly why compactness gives the required uniform tail, suppose
otherwise and choose radii $R_k\to\infty$ and normalized ground
states $\Psi_k$ with
\[
 \int_{|x|>R_k}\rho_{\Psi_k}(x)\,\dd x>\frac12.
\]
After a subsequence $\Psi_k\to\Psi$ strongly in $L^2$.  The bounded
multiplication operator $\sum_i\1_{\{|x_i|>R_k\}}$ has norm at most
$Z$, so its expectations in $\Psi_k$ and $\Psi$ differ by at most
$2Z\|\Psi_k-\Psi\|_2$.  Its expectation in the fixed $\Psi$ tends to
zero by dominated convergence, a contradiction.  Thus a finite
upper radius works uniformly for this $Z$.

For the lower bound, the kinetic estimate
\eqref{eq:atomic-global-density}, with the fixed finite offset $E_Z-E$, gives
\[
 \int\rho_{\Psi_Z}^{5/3}\le C_Z
\]
uniformly over its normalized ground states.  H\"older's inequality
therefore yields
\[
 \int_{|x|<r}\rho_{\Psi_Z}(x)\,\dd x
 \le |B(0,r)|^{2/5}
       \left(\int\rho_{\Psi_Z}^{5/3}\right)^{3/5}
 \le C_Z r^{6/5}.
\]
Choose a positive $r_Z$ making this less than $Z-1/2$.  The exterior
mass then exceeds $1/2$, uniformly over the ground states, so
$R(\Psi_Z)\ge r_Z$.

Taking the minimum of $s_*$ and these finitely many positive lower
radii, and the maximum of $R_*$ and the finitely many upper radii,
gives universal $0<c\le C<\infty$.  All density measures used above
are absolutely continuous.  Their exterior masses are consequently
continuous in the radius, start at $Z$, and tend to zero.  Thus the
infimum defining the half-electron radius has the stated bounds for
every normalized neutral ground state, as required.
\end{proof}
